@@M0@@
Since @@M1@@ is invertible, we have @@M2@@, hence @@M3@@.
Consequently, we have:
@@M4@@

Furthermore, the functoriality of @@M5@@, and the projection formula 3.7 imply:
@@M6@@
Now @@M7@@. Thus:
@@M8@@

Since @@M9@@, @@M10@@ is the image in @@M11@@ of @@M12@@, and is therefore @@M13@@. Furthermore, we have:
@@M14@@

Thus finally we obtain:
@@M15@@
which proves the required relation.

\sgaStableTarget{sga6.e08.sec.5}\section*{5. Riemann--Roch theorem: case of the structural morphism of a projective bundle}

\sgaStableTarget{sga6.e08.para.5-1}\textbf{5.1.} By 3.9, to prove the Riemann--Roch theorem, it remains to prove it for the structural morphism @@M16@@, where @@M17@@ is a quasi-compact scheme having an ample sheaf. Put @@M18@@, @@M19@@, and @@M20@@.
